\documentclass[sigplan,10pt,review,anonymous]{acmart}

\setcopyright{acmcopyright}
\acmConference[PLDI '26]{Programming Language Design and Implementation}{2026}{Worldwide}

%% ============================================================
%%  Packages
%% ============================================================
\usepackage{amsmath,amssymb,amsthm}
\usepackage{mathpartir}          % inferrule macro
\usepackage{listings}
\usepackage{xcolor}
\usepackage{booktabs}
\usepackage{stmaryrd}            % semantic brackets
\usepackage{microtype}

%% ============================================================
%%  Turn syntax highlighting
%% ============================================================
\definecolor{turnblue}{rgb}{0.06,0.39,0.67}
\definecolor{turngreen}{rgb}{0.13,0.55,0.13}
\definecolor{turngray}{rgb}{0.46,0.46,0.46}
\lstdefinelanguage{Turn}{
  keywords={let,return,if,else,while,fn,tool,infer,suspend,spawn,send,receive,
             link,monitor,call,remember,recall,confidence,with,secret,use,
             struct,true,false,null,turn},
  keywordstyle=\color{turnblue}\bfseries,
  comment=[l]{//},
  morecomment=[s]{/*}{*/},
  commentstyle=\color{turngray}\itshape,
  stringstyle=\color{turngreen},
  string=[b]",
  morestring=[b]',
  sensitive=true
}
\lstset{language=Turn, basicstyle=\ttfamily\small, frame=single,
        xleftmargin=1em, numbers=left, numberstyle=\tiny\color{turngray}}

%% ============================================================
%%  Theorem environments
%% ============================================================
\newtheorem{theorem}{Theorem}
\newtheorem{lemma}[theorem]{Lemma}
\newtheorem{definition}[theorem]{Definition}
\newtheorem{corollary}[theorem]{Corollary}

%% ============================================================
%%  Semantic macros
%% ============================================================
\newcommand{\reduces}{\longrightarrow}
\newcommand{\reducesmany}{\longrightarrow^{*}}
\newcommand{\vmstate}[1]{\langle #1 \rangle}
\newcommand{\typecheck}[3]{#1 \vdash #2 : #3}
\newcommand{\infer}{\mathbf{infer}}
\newcommand{\suspend}{\mathbf{suspend}}
\newcommand{\spawn}{\mathbf{spawn}}
\newcommand{\conf}[2]{\mathbf{confidence}(#1,\, #2)}
\newcommand{\uncertain}[2]{\mathbf{Uncertain}(#1,\, #2)}
\newcommand{\result}[1]{\mathbf{Result}\langle #1 \rangle}
\newcommand{\captype}{\mathbf{Cap}}

%% ============================================================
%%  Title block
%% ============================================================
\title{Turn: A Systems Language for Agentic Computation}
\subtitle{Formal Semantics, Distributed Runtime, and Provider-Agnostic Inference}

\author{Muyukani Ephraim Kizito}
\affiliation{\institution{Independent Research}}
\email{muyukaniephraim@yahoo.com}

%%
\begin{document}
\maketitle

%% ============================================================
\begin{abstract}
%% ============================================================
Agentic software—programs that interleave deterministic computation with
probabilistic LLM inference, long-lived state, and distributed message-passing
—imposes constraints that existing general-purpose languages do not model at
the semantic level. Practitioners are forced to layer ad-hoc retry logic,
context management, API key handling, and message routing on top of languages
designed for synchronous, deterministic computation.

We present \textbf{Turn}, a compiled, statically-typed language whose
primitives encode the physical constraints of agentic computation: (1) finite
context capacity, (2) stochastic inference, (3) durable state across
suspension boundaries, and (4) budgeted execution over token expenditure.
Turn's most critical guarantee is \emph{provider agnosticism}: \texttt{infer}
is a first-class language construct that compiles to a generic host trap, not
an HTTP call. LLM provider migrations require no changes to Turn programs.

We describe Turn's type system—including a novel \texttt{Uncertain(T, f64)}
probabilistic type and an object-capability type \texttt{Cap}—its operational
semantics for suspension and resumption, a distributed actor model with
globally-addressed PIDs, and a reference implementation consisting of a
multi-threaded Tokio-backed VM with an HNSW semantic memory index using
Ebbinghaus exponential decay for episodic memory garbage collection.

\end{abstract}

%%
\keywords{programming languages, agentic systems, LLM, operational semantics,
          actor model, object-capabilities, probabilistic types}

%% ============================================================
\section{Introduction}
%% ============================================================

Large language models have created a new category of software: \emph{agents}.
An agent is a long-running program that alternates between deterministic
computation and stochastic inference, maintains state across many interaction
turns, and coordinates work across parallel sub-processes. Building such
systems today requires manually encoding several non-trivial constraints in
application code:

\begin{enumerate}
  \item \emph{Context overflow avoidance} — the model's context window is
        finite; state must be actively managed and summarized.
  \item \emph{Schema enforcement} — LLM outputs are strings; structured
        values must be parsed and retried on schema violation.
  \item \emph{Durability} — agents must survive process restarts without
        replaying the entire inference history.
  \item \emph{Vendor coupling} — API schemas, endpoint paths, and model
        names change; every change can break deployed agents.
  \item \emph{Secret leakage} — API keys are routinely passed through
        prompt templates, exposing them to log pipelines and the model itself.
\end{enumerate}

Turn addresses these constraints at the \emph{language level}, eliminating
the need for framework conventions or ad-hoc implementation patterns.

\paragraph{Contributions.}
\begin{itemize}
  \item A type system with first-class probabilistic values
        ($\uncertain{T}{p}$, $p \in [0,1]$) and an object-capability
        primitive (\captype{}) with enforced non-serializability.
  \item Formal small-step operational semantics for the \texttt{infer},
        \texttt{suspend}, and \texttt{spawn} primitives.
  \item A provider-agnostic inference model in which \texttt{infer}
        compiles to a typed host trap; the language has no knowledge of
        REST endpoints or vendor-specific schemas.
  \item A distributed actor model with globally-addressed PIDs and an
        Erlang-style supervisor tree for fault propagation.
  \item A reference implementation: a multi-threaded Tokio-based bytecode
        VM with an HNSW semantic memory index and Ebbinghaus decay GC.
\end{itemize}

%% ============================================================
\section{Background and Related Work}
%% ============================================================

\paragraph{Agent frameworks.}
LangChain~\cite{langchain}, AutoGen~\cite{autogen}, and CrewAI are library
frameworks for Python. They provide high-level APIs for chaining LLM calls but
do not provide language-level guarantees for type safety, suspension, or secret
isolation. Error recovery and provider migration are programmer responsibilities.

\paragraph{Actor languages.}
Erlang~\cite{erlang} and Elixir provide the actor model that Turn's
\texttt{spawn}/\texttt{send}/\texttt{receive}/\texttt{link} primitives build upon.
Turn extends this model with \emph{distributed PIDs} that are addressable across
physical machines without user-level networking code.

\paragraph{Effect systems.}
Koka~\cite{koka} and Frank~\cite{frank} model computational effects and handlers
as first-class semantic constructs. Turn's \texttt{suspend} and \texttt{infer}
are typed effect boundaries in a similar spirit, though Turn's design prioritizes
agentic-specific constraints (token budgets, semantic memory) over effect generality.

\paragraph{Probabilistic types.}
Languages such as Uncertain\textless T\textgreater~\cite{uncertain} have explored
first-class uncertainty propagation. Turn's \texttt{Uncertain(T, f64)} type is
operationally integrated with \emph{confidence execution traps} that halt
computation when uncertainty violates a declared threshold.

\paragraph{Object-capabilities.}
The object-capability model~\cite{ocap} has been applied in languages such as
E and Pony. Turn applies OCap specifically to the agentic security problem:
LLM tool-call schemas must never include credential values. The \captype{} type
enforces this at the VM level via privilege-violation traps.

%% ============================================================
\section{The Turn Language}
%% ============================================================

\subsection{Syntax Overview}

A Turn program is a sequence of statements. The core grammar is given in
Figure~\ref{fig:grammar}. Expressions include literals, identifiers,
struct initialization, the \texttt{infer} expression, tool invocation via
\texttt{call}, embedding similarity via $\mathtt{\sim>}$, and probabilistic
confidence extraction.

\begin{figure}[t]
\begin{lstlisting}[caption={Core Turn grammar (excerpt)}, label=fig:grammar]
stmt   ::= let x [: Type] = expr ;
         | return expr ;
         | expr ;

expr   ::= literal | id
         | infer Type { expr ; }
         | infer Type with [expr,...] { expr ; }
         | call(expr, expr)
         | spawn(expr) | send pid , expr
         | receive | suspend
         | link(pid) | monitor(pid)
         | remember(expr, expr) | recall(expr)
         | confidence(expr)
         | tool turn(params) -> Type { stmt* }
         | spawn_remote(expr, expr)

Type   ::= Str | Num | Bool | Null | Any
         | List<Type> | Map<Type, Type>
         | Struct(name) | Cap | Uncertain(Type)
         | Result<Type, Type> | Option<Type>
         | Pid | Fn(Type) -> Type
\end{lstlisting}
\end{figure}

\subsection{Type System}

Turn is statically typed with inference. The judgment form is:
$$\typecheck{\Gamma}{e}{\tau}$$
Selected typing rules are given in Figure~\ref{fig:typing}.

\begin{figure}[t]
\begin{mathpar}
\inferrule*[left=T-Infer]
  { \typecheck{\Gamma}{e}{\mathtt{Str}}
  \\ \tau \in \mathit{CognitiveTypes} }
  { \typecheck{\Gamma}{\infer_\tau\{e\}}{\uncertain{\tau}{}} }

\inferrule*[left=T-Uncertain-Elim]
  { \typecheck{\Gamma}{e}{\uncertain{\tau}{}} 
  \\ p \in [0,1] }
  { \typecheck{\Gamma}{\conf{e}{p}}{\mathtt{Bool}} }

\inferrule*[left=T-Cap-Intro]
  { x : \captype{} \in \Gamma }
  { \typecheck{\Gamma}{x}{\captype{}} }

\inferrule*[left=T-Cap-NoSerialize]
  { \typecheck{\Gamma}{e}{\captype{}} }
  { \typecheck{\Gamma}{\mathtt{print}(e)}{\bot}
    \quad \textrm{(PrivilegeViolation)} }

\inferrule*[left=T-Spawn]
  { \typecheck{\Gamma}{e}{\mathtt{Fn}(\tau_1) \to \tau_2} }
  { \typecheck{\Gamma}{\spawn(e)}{\mathtt{Pid}} }

\inferrule*[left=T-Suspend]
  { }
  { \typecheck{\Gamma}{\suspend}{\mathtt{Null}} }
\end{mathpar}
\caption{Selected typing rules.}
\label{fig:typing}
\end{figure}

\paragraph{Probabilistic values.}
The $\uncertain{T}{}$ type represents a value produced by stochastic inference.
It carries a confidence score $p \in [0,1]$. Operators over uncertain values
propagate uncertainty: $\uncertain{T_1}{p_1} + \uncertain{T_2}{p_2}$ yields
$\uncertain{T}{p_1 \cdot p_2}$.

\paragraph{Object-capabilities.}
\captype{} is an opaque integer handle referencing a host-side secret. Rule
T-Cap-NoSerialize makes printing or serializing a capability a static and
dynamic type error respectively. The value never crosses the VM boundary.

%% ============================================================
\section{Operational Semantics}
%% ============================================================

\subsection{VM State}

At any instant, a Turn process state $S$ is a tuple:

$$S = \vmstate{pc,\, \sigma,\, \phi,\, \rho,\, \mathcal{M},\, \mathcal{B}}$$

where $pc$ is the instruction pointer, $\sigma$ is the value stack, $\phi$ is
the call-frame stack, $\rho$ is the environment (scope chain), $\mathcal{M}$
is the semantic memory (HNSW index), and $\mathcal{B}$ is the token-budget
counter.

\subsection{Reduction Rules}

The reduction relation $S \reduces S'$ defines single-step execution.
Selected rules are given in Figure~\ref{fig:semantics}.

\begin{figure}[t]
\begin{mathpar}

\inferrule*[left=E-Infer]
  { \sigma = v \cdot \sigma'
  \\ \textit{hostTrap}(\texttt{llm\_infer}, v) = \uncertain{w}{p} }
  { \vmstate{pc_\mathtt{Infer}, \sigma, \phi, \rho, \mathcal{M}, \mathcal{B}} \reduces
    \vmstate{pc{+}1, \uncertain{w}{p} \cdot \sigma', \phi, \rho, \mathcal{M}, \mathcal{B}} }

\inferrule*[left=E-Suspend]
  { S' = \textit{serialize}(S) }
  { \vmstate{pc_\mathtt{Suspend}, \sigma, \phi, \rho, \mathcal{M}, \mathcal{B}} \reduces
    \textit{Write}(S') }

\inferrule*[left=E-Resume]
  { S = \textit{deserialize}(\textit{Read}()) }
  { \textit{boot} \reduces S }

\inferrule*[left=E-Spawn]
  { \textit{fresh}(\textit{pid}) }
  { \vmstate{pc_\mathtt{Spawn}, cl \cdot \sigma, \phi, \rho, \mathcal{M}, \mathcal{B}} \reduces
    \vmstate{pc{+}1, \textit{pid} \cdot \sigma, \phi, \rho, \mathcal{M}, \mathcal{B}}
  \quad | \quad \textit{fork}(cl, \textit{pid}) }

\inferrule*[left=E-ConfTrap]
  { \uncertain{v}{p} \in \sigma
  \\ p < \theta_S }
  { \vmstate{pc_\mathtt{If}, \sigma, \phi, \rho, \mathcal{M}, \mathcal{B}} \reduces
    \textit{Trap}(\texttt{ConfidenceViolation}) }

\inferrule*[left=E-CapTrap]
  { \captype{} \in \sigma }
  { \vmstate{pc_\mathtt{Print}, \sigma, \phi, \rho, \mathcal{M}, \mathcal{B}} \reduces
    \textit{Trap}(\texttt{PrivilegeViolation}) }

\end{mathpar}
\caption{Small-step operational semantics (selected rules). $\theta_S$ is the
         process strictness threshold; \textit{hostTrap} is the host dispatch interface.}
\label{fig:semantics}
\end{figure}

\subsection{Provider Agnosticism}

Rule E-Infer delegates to $\textit{hostTrap}(\texttt{llm\_infer}, -)$, a
function defined by the \emph{host runtime}, not the language specification.
The Turn AST has no representation of vendor-specific concepts such as REST
endpoints, model names, or proprietary schema formats. Changes to LLM provider
APIs require updating only the host adapter; Turn programs are unaffected.

\subsection{Semantic Memory and Ebbinghaus Decay}

The memory $\mathcal{M}$ is an HNSW approximate nearest-neighbor index. Each
entry stores a value, an embedding vector, and temporal metadata
$(t_{\textit{created}}, t_{\textit{accessed}}, \lambda)$.

Retrieval strength at time $t$ is defined as:
$$\mathit{RS}(q, m, t) = \cos(\vec{q}, \vec{m}) \cdot e^{-\lambda (t - t_{\textit{accessed}})}$$

A background garbage-collector prunes entries whose $\mathit{RS}$ falls below
a global noise threshold $\varepsilon$. Direct recall resets
$t_{\textit{accessed}} \coloneq \textit{now}$, implementing spaced repetition.

%% ============================================================
\section{Distributed Actor Model}
%% ============================================================

Processes in Turn are globally addressable via structured PIDs:
$$\textit{Pid} ::= \{ \textit{node\_id} : \mathtt{Str},\; \textit{local\_id} : \mathbb{N} \}$$

A \emph{Switchboard} running on each node routes messages. If
$\textit{pid}.\textit{node\_id} = \texttt{"local"}$, the message is deposited
into an in-process \texttt{tokio::mpsc} channel. Otherwise, the Switchboard
serializes the message and dispatches it over TCP/gRPC to the target node's
Switchboard.

\texttt{spawn\_remote("node\_ip", closure)} deploys a closure onto a remote
node. Object-capabilities held locally cannot be serialized over the wire; the
Capability Routing Protocol proxies tool execution back to the originating
host's registry to satisfy remote \texttt{infer} calls that require credential
access.

\paragraph{Supervisor Trees.}
\texttt{link(pid)} establishes a bi-directional link between two processes.
When a process terminates (normally or due to an unhandled error), the VM
emits a $\textsc{ProcessExit}$ signal to all linked PIDs. The receiving process
may trap the signal to implement restart strategies, mirroring Erlang's OTP
supervision model~\cite{erlang}.

%% ============================================================
\section{Reference Implementation}
%% ============================================================

\subsection{Architecture}

The Turn v0.5.0-alpha reference implementation is a Rust bytecode VM. The
pipeline is: \textbf{Lexer} $\to$ \textbf{Recursive-Descent Parser} $\to$
\textbf{AST} $\to$ \textbf{Semantic Analyzer} $\to$ \textbf{Bytecode Compiler}
$\to$ \textbf{Async VM}.

\subsection{Async Scheduler}

Each spawned Turn process maps to a \texttt{tokio::spawn} green thread. The
VM executor is \texttt{async fn}, enabling cooperative multi-tasking without
blocking the thread pool. Host traps (LLM calls, tool execution) suspend the
current task and resume it upon completion, allowing other processes to run.

\subsection{Implementation Status}

\begin{table}[h]
\centering
\begin{tabular}{@{}llr@{}}
\toprule
Subsystem & Status & Notes \\
\midrule
Lexer / Parser / AST      & \checkmark & Full Turn grammar \\
Bytecode Compiler         & \checkmark & Type-annotated instruction set \\
Async VM Executor         & \checkmark & Tokio work-stealing \\
\texttt{infer} + Cognitive Types & \checkmark & Provider-agnostic host trap \\
\texttt{suspend} / WAL    & \checkmark & Cross-restart replay \\
Actor Model               & \checkmark & Isolated tokio::mpsc mailboxes \\
Supervisor Trees          & \checkmark & Erlang-style fault propagation \\
HNSW Semantic Memory      & \checkmark & Ebbinghaus decay GC \\
Arc Structural Sharing    & \checkmark & Zero-copy message passing \\
Object-Capability (Cap)   & \checkmark & PrivilegeViolation traps \\
Generic Types (List\textless T\textgreater) & \checkmark & AST + schema generator \\
Monadic Errors            & \checkmark & Result\textless T,E\textgreater, Option\textless T\textgreater \\
Confidence Traps          & \checkmark & Strictness threshold enforced \\
Distributed Swarms        & \checkmark & TCP/gRPC Switchboard \\
Standard Library          & \checkmark & fs, http, math, json, time, regex \\
Language Server (LSP)     & \checkmark & Hover, go-to-definition, diagnostics \\
VS Code Extension         & \checkmark & Syntax + LSP client \\
Provider Adapters         & \checkmark & OpenAI, Anthropic, Gemini, Ollama \\
\bottomrule
\end{tabular}
\caption{Implementation status as of v0.5.0-alpha.}
\end{table}

%% ============================================================
\section{Example Programs}
%% ============================================================

\subsection{Provider-Agnostic Inference with Type Safety}

\begin{lstlisting}
struct Invoice { amount: Num, vendor: Str }

let invoice = infer Invoice {
    "Extract invoice details from: " + raw_text;
};

// The VM retried on schema mismatch.
// `invoice` is guaranteed to be a valid Invoice.
print(invoice.amount);
\end{lstlisting}

\subsection{Secure Tool with Capability}

\begin{lstlisting}
let payment_tool = tool turn(
    amount: Num,
    secret api_key: Cap    // never in LLM schema
) -> Str {
    return call("stripe_charge", amount);
};

let result = infer Str with [payment_tool] {
    "Charge the customer $42 for invoice #1234";
};
\end{lstlisting}

\subsection{Distributed Supervisor}

\begin{lstlisting}
let worker_pid = spawn_remote("10.0.0.5", fn() {
    while true {
        let job = receive;
        let result = infer Str { job; };
        send supervisor_pid, result;
    }
});

link(worker_pid);
// If worker crashes, this process receives a
// ProcessExit signal and may restart it.
\end{lstlisting}

%% ============================================================
\section{Discussion}
%% ============================================================

\paragraph{What Turn is not.}
Turn is not a general-purpose systems language competing with Rust or C++.
Its primitive set is deliberately narrow and agentic-specific. A Turn program
is expected to call out to tools and external services for computation outside
its cognitive domain.

\paragraph{Token budgets.}
The current reference implementation tracks instruction cycles and delegates
token counting to host adapters (since tokenization is model-specific). Future
work will standardize a token-budget opcode that the VM enforces before each
\texttt{infer} dispatch.

\paragraph{Formal verification.}
The type safety theorem (progress + preservation) for Turn's core calculus is
left as formal future work. The operational semantics presented here provide
the necessary foundation; we expect progress and preservation to hold by
construction of the host-trap boundary.

%% ============================================================
\section{Conclusion}
%% ============================================================

Turn demonstrates that the constraints of agentic computation—stochastic
inference, durable state, finite context capacity, secret isolation, and
distributed coordination—can be expressed as first-class language primitives
with sound operational semantics. The result is a language where programs are
immune to LLM provider churn, agent secrets cannot leak through LLM schemas,
and distributed actor orchestration requires no external middleware.

The v0.5.0-alpha reference implementation delivers all described features and
is available at \url{https://github.com/ekizito96/Turn}.

%% ============================================================
\bibliographystyle{ACM-Reference-Format}
\begin{thebibliography}{9}

\bibitem{langchain}
Chase, H. (2022). \textit{LangChain}. GitHub.
\url{https://github.com/langchain-ai/langchain}

\bibitem{autogen}
Wu, Q. et al. (2023). AutoGen: Enabling Next-Gen LLM Applications via
Multi-Agent Conversation Framework. \textit{arXiv:2308.08155}.

\bibitem{erlang}
Armstrong, J., Virding, R., Wikstr{\"o}m, C., and Williams, M. (1996).
\textit{Concurrent Programming in Erlang}, 2nd ed. Prentice Hall.

\bibitem{koka}
Leijen, D. (2014). Koka: Programming with Row Polymorphic Effect Types.
\textit{MSFP 2014}.

\bibitem{frank}
Lindley, S., McBride, C., and McLaughlin, C. (2017). Do Be Do Be Do.
\textit{POPL 2017}.

\bibitem{uncertain}
Bornholt, J., Mytkowicz, T., and McKinley, K.S. (2014).
Uncertain\textless T\textgreater: A First-Order Type for Uncertain Data.
\textit{ASPLOS 2014}.

\bibitem{ocap}
Miller, M.S. (2006). \textit{Robust Composition: Towards a Unified Approach
to Access Control and Concurrency Control}. PhD thesis, Johns Hopkins University.

\end{thebibliography}

\end{document}
